\chapter{Results}
\label{chap:results}

% TODO: bridging paragraph: all experiments reported against identical
%       splits; headline comparison table appears in Section \ref{sec:overall}.

\section{Data Exploration Findings (Exp 0)}
% TODO: workspace plots, occlusion statistics, Dataset 2 time-series confirmation.

\section{PCC Analytical Baseline (Exp 1)}

\subsection{Single-Segment PCC on Dataset 2 (Exp 1A)}
% TODO: error distribution, failure modes.

\subsection{Two-Segment PCC on Dataset 1 (Exp 1B, if attempted)}
% TODO: note whether the stretch goal was pursued.

\section{MLP IK (Exp 2)}

\subsection{Tip-Only MLP (Exp 2A)}
% TODO: results on Dataset 1 and Dataset 2.

\subsection{Tip + Mid Ablation on Dataset 1 (Exp 2B)}
% TODO: does mid-marker input reduce error? Expected: yes, due to
%       two-segment underdetermination.

\section{Temporal LSTM IK (Exp 3)}
% TODO: comparison against MLP; hysteresis analysis.

\section{Jacobian-Learned IK (Exp 4)}
% TODO: convergence behaviour, final tip error, robustness.

\section{Sim-to-Real with PyElastica (Exp 5)}
% TODO: sim-only baseline vs. fine-tuned on real Dataset 2.

\section{Trajectory Tracking (Exp 6)}
% TODO: uniform protocol, path-following error, actuation smoothness.

\section{Overall Comparison}
\label{sec:overall}
% TODO: master table (MAE / RMSE / MAPE per method per dataset) and a
%       one-paragraph commentary on the best-performing method.

\section{Ablation Studies}
% TODO: mid-marker ablation, sequence length, damping factor for DLS,
%       simulator fine-tune ratio.
